% !TEX program = lualatex
\documentclass[12pt, a4paper]{article}
\usepackage{master}

\DeclareWorkGR{Cat}{범주론}{Κατηγορίαι}{Categoriae}
\DeclareWorkGR{DI}{명제론}{Περὶ ἑρμηνείας}{De interpretatione}
\DeclareWorkGR{Phys}{자연학}{Φυσικά}{Physica}
\DeclareWorkGR{Met}{형이상학}{Μεταφυσικά}{Metaphysica}

\fancyhead[L]{\footnotesize\color{midgray}오르가논 강독}
\fancyhead[R]{\footnotesize\color{midgray}제5주}

\begin{document}
\thispagestyle{firstpage}

\weeksection{제5주}{범주론 11a38--15b32}{나머지 범주와 후범주론}

\subsection{나머지 범주의 약론 (9장, 11b1--14)}

범주론의 6--8장이 양·관계·질을 상세히 다루는 것과 대조적으로, 9장은 나머지 여섯 범주 --- 행위(\gr{ποιεῖν}), 겪음(\gr{πάσχειν}), 장소(\gr{ποῦ}), 시간(\gr{ποτέ}), 상태(\gr{κεῖσθαι}), 소유(\gr{ἔχειν}) --- 를 극도로 간략하게 처리합니다. 아리스토텔레스는 이것들이 `명백하기 때문에'(11b14) 더 이상 말할 것이 없다고 합니다. 실제로 9장에서 명시적으로 논의되는 것은 행위와 겪음뿐이며, 그마저도 이 범주들이 반대를 허용하고 정도의 차이를 허용한다는 것이 전부입니다.

이 간략함이 진정으로 `명백하기 때문'인지, 아니면 아리스토텔레스가 이 범주들에 대한 상세한 분석을 아직 완성하지 못했기 때문인지는 논쟁의 여지가 있습니다. 올레(Oele)는 행위와 겪음이 학계에서 `특히 소홀히 다루어져' 왔다고 지적하면서, 이 두 범주가 \citework{Phys}{III.3}, 생성소멸론, 영혼론에서는 핵심적인 역할을 한다고 강조합니다.\footnote{Marjolein Oele, ``Understanding Aristotle's Poiein and Paschein,'' University of San Francisco, 2012. \citework{Phys}{III.3}에서 아리스토텔레스는 행위와 겪음이 운동 자체와 존재론적으로 구별되지 않으며, 같은 운동을 행위자의 관점에서 보면 행위이고 수용자의 관점에서 보면 겪음이라고 설명합니다.}

\subsection{후범주론: 논리적 어휘의 정비 (10--15장, 11b15--15b32)}

\subsubsection{후범주론의 성격과 진정성 문제}

10장부터 15장까지의 내용은 전통적으로 `후범주론'(Postpraedicamenta)이라 불립니다. 이 부분은 범주 자체를 다루는 것이 아니라, 범주 분석에 사용되는 논리적 개념들 --- 대립(\gr{ἀντίκεισθαι}), 선후(\gr{πρότερον}), 동시(\gr{ἅμα}), 변화(\gr{κίνησις}), 소유(\gr{ἔχειν}) --- 를 정의하고 분류합니다.

후범주론의 진정성(authenticity) 문제는 고대부터 제기되어 왔습니다. 프레데(Frede)에 따르면, `특히 후범주론에 대한 의심은 결코 완전히 사라진 적이 없습니다.'\footnote{Michael Frede, ``The Title, Unity, and Authenticity of the Aristotelian Categories,'' in \gri{Essays in Ancient Philosophy} (Minneapolis: University of Minnesota Press, 1987), 11--28.} 로스(Ross)는 후범주론이 `일반적으로 위작으로 간주된다'고 말했습니다. 반면 애크릴(Ackrill)은 본론과 후범주론을 연결하는 이행 문장(11b10--16)이 아리스토텔레스의 것이 아니라고 보면서도, 후범주론의 내용 자체의 진정성까지 부정하지는 않았습니다.\footnote{Ackrill, \gri{Aristotle's Categories and De Interpretatione}, 70.}

\subsubsection{네 가지 대립 (10장, 11b17--13b35)}

10장은 후범주론에서 가장 중요한 장이며, 이후 오르가논 전체의 논리적 어휘를 제공합니다.

\srcquote{\gr{Ἀντικεῖσθαι δὲ ἕτερον ἑτέρῳ τετραχῶς λέγεται, ἢ ὡς τὰ πρός τι, ἢ ὡς τὰ ἐναντία, ἢ ὡς στέρησις καὶ ἕξις, ἢ ὡς κατάφασις καὶ ἀπόφασις.}}{`하나가 다른 하나에 대립한다는 것은 네 가지 방식으로 말해진다: 관계로서, 반대로서, 결여와 소유로서, 긍정과 부정으로서.'}{범주론 11b17--19}

{%
\footnotesize
\setlength{\extrarowheight}{1.5pt}%
\renewcommand{\arraystretch}{1.1}%
\setlength{\tabcolsep}{3pt}%
\rowcolors{2}{altrowbg}{white}%
\noindent\centering
\begin{tabular}{|>{\centering\arraybackslash}m{18mm}|>{\centering\arraybackslash}m{22mm}|>{\centering\arraybackslash}m{28mm}|>{\centering\arraybackslash}m{18mm}|>{\centering\arraybackslash}m{22mm}|}
\hline
\rowcolor{tableheadblue}
{\color{h1navy}\bfseries 유형} &
{\color{h1navy}\bfseries 그리스어} &
{\color{h1navy}\bfseries 예시} &
{\color{h1navy}\bfseries 상호 정의} &
{\color{h1navy}\bfseries 참·거짓\newline 필연 분배} \\ \hline
관계적 & \gr{πρός τι} & 두 배 / 절반 & 예 & 아니오 \\ \hline
반대 & \gr{ἐναντία} & 좋음 / 나쁨 & 아니오 & 아니오 \\ \hline
결여·소유 & \gr{στέρησις}/\gr{ἕξις} & 눈멀음 / 시력 & 아니오 & 아니오 \\ \hline
긍정·부정 & \gr{κατάφασις}/\gr{ἀπόφασις} & 앉아 있다 / 앉아 있지 않다 & 아니오 & 예 \\ \hline
\end{tabular}
\par
}

이 네 가지 가운데 논리학적으로 가장 중요한 것은 긍정과 부정의 대립입니다. 아리스토텔레스는 이것만이 반드시 하나가 참이고 하나가 거짓이라는 특성을 갖는다고 주장합니다(13a37--b26). 반대의 경우에는 둘 다 거짓일 수 있습니다 --- 어떤 것이 좋지도 나쁘지도 않을 수 있습니다. 이 구별은 명제론에서 모순(\gr{ἀντίφασις})과 반대(\gr{ἐναντίωσις})의 구별로 발전하며, 중세 논리학의 `대당 사각형'(Square of Opposition)의 기초를 제공합니다.\footnote{대당 사각형 도식 자체는 아리스토텔레스의 것이 아니라 아풀레이우스(Apuleius, 2세기)에게서 처음 나타나고 보에티우스에 의해 유명해졌습니다.}

\subsubsection{반대와 매개항 (11장, 13b36--14a25)}

11장은 반대를 더 세분합니다. 어떤 반대 쌍에는 매개항(\gr{μεταξύ})이 있고, 어떤 반대 쌍에는 없습니다. 검정과 흼 사이에는 회색이라는 매개항이 있습니다 --- 어떤 것이 검지도 희지도 않을 수 있습니다. 반면 홀수와 짝수 사이에는 매개항이 없습니다 --- 수는 반드시 홀수이거나 짝수입니다.

\subsubsection{선후의 다섯 가지 의미 (12장, 14a26--14b23)}

\srcquote{\gr{Πρότερον ἕτερον ἑτέρου λέγεται τετραχῶς. πρῶτον μὲν καὶ κυριώτατα κατὰ χρόνον, καθ᾽ ὃ πρεσβύτερον ἕτερον ἑτέρου λέγεται καὶ παλαιότερον.}}{`하나가 다른 것보다 선행한다는 것은 네 가지로 말해진다. 첫째이자 가장 본래적인 것은 시간에 따른 것으로, 그에 따라 하나가 다른 것보다 더 나이 들었다거나 더 오래되었다고 말해진다.'}{범주론 14a26--27}

아리스토텔레스는 네 가지를 열거한 뒤 다섯 번째를 추가합니다. 시간적 선후, 비대칭적 존재 함축(1이 존재하면 2가 존재하지만 역은 아니다), 어떤 순서에 따른 선후(전제는 결론에 선행), 가치에 따른 선후(더 나은 것이 선행), 자연적 선후(원인이 되는 쪽이 선행)입니다.\footnote{Ana Laura Edelhoff, \gri{Aristotle on Ontological Priority in the Categories}, Cambridge Elements in Ancient Philosophy (Cambridge: Cambridge University Press, 2020).}

\subsubsection{동시의 세 가지 의미 (13장)}

13장은 `동시'(\gr{ἅμα})의 세 가지 의미를 분석합니다. 시간적 동시(같은 시간에 생겨남), 자연적 동시(존재의 측면에서 서로 함축하되 어느 쪽도 원인이 아닌 것 --- 두 배와 절반), 같은 류의 같은 분할에서 나온 종들의 동시(새와 물고기와 육상동물)입니다.

\subsubsection{변화의 여섯 가지 (14장, 15a13--15b16)}

{%
\footnotesize
\setlength{\extrarowheight}{1.5pt}%
\renewcommand{\arraystretch}{1.1}%
\setlength{\tabcolsep}{3pt}%
\rowcolors{2}{altrowbg}{white}%
\noindent\centering
\begin{tabular}{|>{\centering\arraybackslash}m{18mm}|>{\centering\arraybackslash}m{22mm}|>{\centering\arraybackslash}m{28mm}|>{\centering\arraybackslash}m{36mm}|}
\hline
\rowcolor{tableheadblue}
{\color{h1navy}\bfseries 종류} &
{\color{h1navy}\bfseries 그리스어} &
{\color{h1navy}\bfseries 예시} &
{\color{h1navy}\bfseries 관련 범주} \\ \hline
생성 & \gr{γένεσις} & 탄생 & 실체 \\ \hline
소멸 & \gr{φθορά} & 죽음 & 실체 \\ \hline
증가 & \gr{αὔξησις} & 성장 & 양 \\ \hline
감소 & \gr{φθίσις} & 쇠퇴 & 양 \\ \hline
성질 변화 & \gr{ἀλλοίωσις} & 색의 변화 & 질 \\ \hline
장소 이동 & \gr{κατὰ τόπον μεταβολή} & 걷기 & 장소 \\ \hline
\end{tabular}
\par
}

이 분류는 \citework{Phys}{V권}에서 더 상세하게 전개됩니다 --- 자연학에서 아리스토텔레스는 생성과 소멸을 엄밀한 의미의 운동(\gr{κίνησις})에서 제외하고, 운동을 양·질·장소의 세 범주에서의 변화로 한정합니다.

15장은 `소유'(\gr{ἔχειν})의 여러 의미를 열거합니다: 상태를 소유함, 양을 소유함, 옷을 걸침, 부분을 소유함, 담고 있음, 재산을 소유함 등. 이것은 체계적인 분류라기보다는 용법의 열거에 가깝습니다.

\subsubsection{후범주론의 의의}

후범주론이 진정한 아리스토텔레스의 것이든 후대 편집의 산물이든, 그 내용이 범주론의 본론과 긴밀하게 연결되어 있다는 점은 부정하기 어렵습니다. 실체에 반대가 없다는 주장(5장), 양에 반대가 없다는 주장(6장), 관계가 본성상 동시적이라는 주장(7장)은 모두 10장과 12--13장의 개념들을 전제합니다. 후범주론은 이 전제들을 명시적으로 정의함으로써, 본론의 주장들을 정밀하게 진술할 수 있는 개념적 문법을 제공합니다.

다음 주(6주차)부터 우리는 범주론의 틀 위에 세워진 포르피리오스의 이사고게를 읽습니다. 이사고게는 범주론을 전제하면서 거기에 다섯 가지 술어가능자(류·종·종차·고유성·부수성)의 체계를 덧붙이는 작업입니다.

% ── 참고문헌 ──
\subsection*{참고문헌}
\begin{references}

\refitem Ackrill, J.\ L. \gri{Aristotle's Categories and De Interpretatione}. Clarendon Aristotle Series. Oxford: Clarendon Press, 1963.

\refitem Edelhoff, Ana Laura. \gri{Aristotle on Ontological Priority in the Categories}. Cambridge Elements in Ancient Philosophy. Cambridge: Cambridge University Press, 2020.

\refitem Frede, Michael. ``The Title, Unity, and Authenticity of the Aristotelian Categories.'' In \gri{Essays in Ancient Philosophy}, 11--28. Minneapolis: University of Minnesota Press, 1987.

\refitem Minio-Paluello, L., ed. \gri{Aristotelis Categoriae et Liber de Interpretatione}. OCT. Oxford: Clarendon Press, 1949; repr.\ with corrections, 1956. Eng.\ trans.: J.\ L.\ Ackrill, \gri{Aristotle's Categories and De Interpretatione}, Clarendon Aristotle Series. Oxford: Clarendon Press, 1963.

\refitem Oele, Marjolein. ``Understanding Aristotle's Poiein and Paschein.'' University of San Francisco, 2012.

\refitem Studtmann, Paul. ``Aristotle's Categories.'' In \gri{Stanford Encyclopedia of Philosophy}.\newline https://plato.stanford.edu/entries/aristotle-categories/.

\end{references}

\end{document}
